\begin{pastexam}[label={px:ncu-107-dsa-22}]{中央資工所 107 學年度 第 22 題}
（複選；每題 5 分，答錯每小題倒扣 1 分，扣到零分為止）

Below is the algorithm AllPairCost which computes the shortest distances between all pairs of
vertices i, j, where i $\ne$ j. Formally, given distances of edges in graph G, determine the
shortest distances between all pairs of vertices in G. Note that the distance of an arbitrary edge
in G is non-negative.
\begin{quote}
\textbf{Algorithm} AllPairCost\\
\textbf{Input:} a two dimensional array C, where C[i][j]$>$0 denotes the distance of directed edge
(i, j) (i.e., the edge from vertex i to vertex j).\\
\textbf{Output:} a two dimensional array D, where D[i][j] denotes the shortest distance from
vertex i to vertex j.
\end{quote}
\begin{minted}[fontsize=\small,breaklines]{c}
1:  int i, j, k;
2:  for ( i=0; i<n; i++ )
3:     for ( j=0; j<n; j++ )
4:          D[i][j] = C[i][j];
5:  for ( k=0; k<n; k++ )
6:     for ( i=0; i<n; j++ )
7:          for ( j=0; j<n; j++ )
8:              if ( ___(B7)___ )
9:                    ____(B8)____;
\end{minted}

Which of the following statements is true?
\begin{enumerate}[label=(\alph*)]
\item For two non-adjacent vertices u and v, C[u][v] should be 0
\item Algorithm AllPairCost is applicable to graphs with cycles
\item Algorithm AllPairCost is applicable to undirected graphs
\item Algorithm AllPairCost is applicable to graphs with multiple edges (i.e., two or more edges
are incident to the same two vertices)
\item All of the above
\end{enumerate}
\end{pastexam}

\begin{pastsoln}
答案：(b)(c)(d)。

AllPairCost 是 Floyd--Warshall（歷屆試題~\ref{px:ncu-107-dsa-20}，(B7)、(B8) 填 (b)，不等號應為 \texttt{>}）。
題目說邊的距離非負，所以圖上沒有負環。

\begin{itemize}
\item (a) \textbf{錯}。 不相鄰的 $u$、$v$ 之間沒有邊，$C[u][v]$ 應該設成 $\infty$（實作上用一個比任何路徑都
大的數），表示「目前走不到」（資結書第~25.6~節的 $D^{(0)}$）。 設成 $0$ 等於說 $u$ 到 $v$ 有一條長度 $0$ 的
邊，最短距離會全部算錯。 例如兩個頂點、只有一條邊 $0 \to 1$ 長 $5$，$C[1][0]$ 設成 $0$ 就會得到
$D[1][0] = 0$，但 $1$ 根本走不到 $0$。 （對角線 $C[i][i]$ 才是 $0$。）
\item (b) \textbf{對}。 Floyd--Warshall 只要求沒有負環。 邊的距離非負，任何環的總長都 $\ge 0$，有環的圖照樣
適用。 （有正環時，最短路徑不會繞環，因為繞一圈只會更長。）
\item (c) \textbf{對}。 無向邊 $\{u, v\}$ 長 $w$，當成兩條有向邊 $u \to v$、$v \to u$，令
$C[u][v] = C[v][u] = w$（相鄰矩陣對稱，資結書第~21.2~節）。 權重非負，來回走一條邊得到的環長 $2w \ge 0$，
不會有負環，演算法適用，算出來的 $D$ 也是對稱的。
\item (d) \textbf{對}。 $C$ 每一對頂點只能存一個數，但同方向的平行邊裡，最短路徑只會用最短的那一條（換成較長
的只會變長），所以輸入時令 $C[u][v]$ 為所有 $u \to v$ 邊裡最短的長度，演算法就照樣適用。
\item (e) \textbf{錯}。 (a) 錯，所以「以上皆是」不成立。
\end{itemize}
題目寫「Which of the following statements is true?」，但它在多選題的區段，(b)(c)(d) 三項都對，全選。
(d) 要先把平行邊取最小值才放得進 $C$，如果認為「放不進 $C$ 就不算適用」，會只選 (b)(c)；本書的看法是這個前處理
是輸入格式的一部分，演算法本身不必改，所以選 (d)。
\end{pastsoln}
